% Section 6.9, from lectures/L06/appendix/b_exercises.md.

\section{Exercises}
\label{c6:sec:exercises}

\begin{aside}
\textbf{How to check your work.} Every exercise below that asks you to write a VHDL module has a
self-checking testbench under \file{lectures/L06/exercises/}. Write your module in its directory
\file{exercises/<module>/}, with the entity name and the \textbf{port order} the exercise states,
and then run it with GHDL \textemdash{} see \secref{c2:sec:testbenches} for the three commands.
\end{aside}

% ------------------------------------------------------------------------------------------
\exerciseset{Counters}

\exercise{Width and wraparound point}{Reasoning}
\label{c6:ex:width}
A signal is declared as:

\begin{vhdlcode}
signal counter: natural range 0 to 1023;
\end{vhdlcode}

\xpart{a} How many bits wide is the register this synthesises to?

\xpart{b} What is the largest value \code{counter} can hold before it wraps back to \code{0} at the
next increment?

\xpart{c} Rewrite the declaration for a counter that instead has to count from \code{0} up to
\code{63}.

Follow the counter pattern in \secref{c6:sec:counter}.

\exercise{Trace a wraparound by hand}{Reasoning}
\label{c6:ex:trace}
Trace a counter of type \code{natural range 0 to 3} by hand (no VHDL needed). Assume the counter
starts at \code{2}, that it is incremented once per clock edge for five consecutive edges, and that
there is no reset in between.

Note, clock edge by clock edge: the value held before the edge, and the value captured at the edge.
Explain why the counter returns to \code{0} with no explicit ``if the count reaches its maximum,
clear it'' logic.

\note[Note]This is the \emph{synthesised} behaviour, the one that matters for hardware.
\Secref{c6:sec:counter} flags the simulation caveat that goes with it: GHDL range-checks a signal of
type \code{natural range 0 to 3}, so a simulation of this counter aborts at the wraparound instead
of rolling over. Trace it on paper as hardware, and reach for a full-width \code{unsigned} if you
ever need a counter that wraps in simulation too.

\exercise{The counter by hand, and the bit that runs out}{Circuit}
\label{c6:ex:counterhand}
Build the 4-bit counter from \secref{c6:sec:counter} by hand in CircuitVerse, and use it to see the
overflow rather than take it on trust.

\xpart{a} Build it: a 4-bit register (four D flip-flops sharing a clock), an adder with the
register's outputs on one input and a constant \code{1} on the other, the adder's sum wired back to
the register's \code{D} inputs, and four output elements or LEDs, one per bit.

\xpart{b} Run it with a clock period of \code{1000 ms} and note what you see:
\begin{itemize}
  \item Write down the sequence of counts over sixteen consecutive clock edges, starting from
    \code{0000}.
  \item Keep an eye on \code{carry_out} on the topmost adder as the count rolls from \code{1111} to
    \code{0000}: during how many of the sixteen edges is it high, and what is connected to it?
  \item Explain, in one sentence, where the fifth bit goes, and why the register stores \code{0000}
    rather than \code{10000}.
\end{itemize}

\xpart{c} Now reason about what you did \emph{not} build: point at the part of your drawing that
decides the counter should return to \code{0000}. There is none. Say what does the job instead. What
would you have to add to make it count \code{0} to \code{9} and then wrap?
(Exercise~\ref{c6:ex:counter} asks you to write exactly that in VHDL, so keep your answer.)

\xpart{d} Change the width to three bits by removing the top flip-flop and narrowing the adder.
Predict the new wraparound point before you run it, confirm it, and state the rule connecting a
counter's width to the value it wraps at.

\textbf{Keep this circuit.} \Chapref{c7:ch} builds its timer by adding two things to it, a
comparator and a clear, and starting from a counter you have already seen working is the whole point
of the exercise there.

\exercise{\code{counter}: a counter with an arbitrary radix}{VHDL}
\label{c6:ex:counter}
Write an entity named \code{counter} that counts from \code{0} up to one below its radix, and then
wraps back to \code{0}. With the default radix that is a decimal counter, \code{0} to \code{9}.

The entity has one generic:

\begin{booktable}{@{}p{4.4em}p{9.8em}p{4em}L@{}}
  \thead{Generic} & \thead{Type} & \thead{Default} & \thead{Meaning} \\
  \midrule
  \code{RADIX} & \code{natural range 1 to 2**16} & \code{10} & How many values the counter goes
    through before repeating. It counts \code{0} to \code{RADIX-1} inclusive, so \code{RADIX} is the
    length of the cycle, not the largest count. \\
\end{booktable}

and these ports:

\begin{booktable}{@{}p{5.6em}p{3.4em}p{10em}L@{}}
  \thead{Port} & \thead{Dir.} & \thead{Type} & \thead{Meaning} \\
  \midrule
  \code{clock} & in & \code{std_logic} & System clock. \\
  \code{reset_s2_n} & in & \code{std_logic} & Active-low, \textbf{already synchronised}: it comes
    from your \code{reset_sync}, never straight from a pin. Asynchronous, so it clears with no clock
    edge involved. \\
  \code{count} & out & \code{natural range 0 to RADIX-1} & The running count. \code{0} is a value it
    takes, so the range starts there: reset drives it to \code{0}, and it returns there at every
    wraparound. \\
  \code{tick} & out & \code{std_logic} & A one-cycle pulse at the edge where \code{count} wraps from
    \code{RADIX-1} back to \code{0}. \\
\end{booktable}

Implement it with a single synchronous process:
\begin{itemize}
  \item Include \code{clock} and \code{reset_s2_n} in the sensitivity list.
  \item At reset, set \code{count} to \code{0} and \code{tick} to \code{'0'}.
  \item At every rising clock edge: drive \code{tick} low by default; when \code{count} has reached
    \code{RADIX-1}, wrap \code{count} to \code{0} and pulse \code{tick} high; otherwise increment
    \code{count}.
\end{itemize}

Write the comparison against \code{RADIX-1}, never against \code{9}. The generic is the whole
difference between a module and a module you can use twice: \code{RADIX} is fixed before synthesis
runs, so the comparator it builds is no larger than a hardcoded one, and the counter's register is
sized for the radix you asked for. That is the same argument \secref{c4:sec:generics} makes for
\code{button_sync}'s \code{COUNT}, and the reason \chapref{c7:ch}'s \code{timer} takes its target
value as a generic rather than as a constant.

Unlike the free wraparound of a counter of type \code{natural range 0 to 15}, a counter with an
arbitrary radix needs an explicit comparison, since \code{RADIX-1} is in general not a power-of-two
boundary. Note what happens when it is: with \code{RADIX = 16} the comparison is still written out,
but the hardware underneath no longer needs it.

\note[Hint]The process needs to read the running count (to compare it and to increment it), so hold
the count in an internal \code{signal count_s: natural range 0 to RADIX-1;}, run the logic on
\code{count_s}, and drive the output with \code{count <= count_s;} outside the process, exactly as
\file{d_flip_flop.vhd} uses \code{q_s} in \chapref{c3:ch}. That is no stylistic choice: every
\code{ghdl} command in this book passes \code{--std=93}, and under VHDL-93 reading an \code{out}
port does not compile at all. Later standards relax that, which is why you may have seen it done
elsewhere.

Follow the counter pattern in \secref{c6:sec:counter}, and the drive-low-then-pulse shape
\code{data_ready} has in \secref{c6:sec:rx8}.

\modulefig[0.68]{lectures/L06/appendix/images/counter.png}

\begin{selfcheck}
\textbf{Self-check.} Name your entity \code{counter}, with the generic \code{RADIX}
(\code{natural range 1 to 2**16}, default \code{10}), the inputs \code{clock}, \code{reset_s2_n} and
the outputs \code{count} (\code{natural range 0 to RADIX-1}) and \code{tick} (\code{std_logic}),
declared in that order; its testbench is in \file{exercises/counter/}. It elaborates \textbf{two}
instances, one with \code{RADIX = 10} and one with \code{RADIX = 4}, and runs each through two full
cycles, checking that \code{tick} is a one-cycle pulse rather than a level and that it never fires
halfway through. An architecture that silently hardcodes \code{10} passes the first instance and
fails on the second, which is the point of checking two.
\end{selfcheck}

% ------------------------------------------------------------------------------------------
\exerciseset{Shift registers}

\exercise{Four bits into a SIPO register}{Reasoning}
\label{c6:ex:sipotrace}
An 8-bit SIPO shift register starts at \code{00000000} and uses the idiom from
\secref{c6:sec:sipo}:

\begin{vhdlcode}
shift_reg <= shift_reg(6 downto 0) & serial_in;
\end{vhdlcode}

It receives the serial bit stream \code{1, 0, 1, 1} (one bit per clock edge, in that order) over
four consecutive rising clock edges.

\xpart{a} Write out the whole 8-bit contents of \code{shift_reg} after each of the four clock edges.

\xpart{b} After the four edges: which four bits of \code{shift_reg} hold the received data, and in
what order do they sit, relative to the order the bits arrived in?

\exercise{Four bits out of a PISO register}{Reasoning}
\label{c6:ex:pisotrace}
A 4-bit PISO register, following the idiom from \secref{c6:sec:piso}, is parallel-loaded with
\code{parallel_in = "1010"} (\code{load = '1'} for one clock edge), then shifted with
\code{shift = '1'} for the following four clock edges, with \code{serial_in} held at \code{'0'}
throughout.

Remember the direction from \secref{c6:sec:shift}: bits move towards the MSB, \code{serial_out} is
bit 3 (the value about to leave), and \code{serial_in} enters at bit 0.

\xpart{a} Write out the register's contents and \code{serial_out} at every step, one row per clock
edge: the row for the load edge, before any shifting, and then one row for each of the four shift
edges.

\xpart{b} In what order do the four loaded bits appear on \code{serial_out}, relative to their
positions in \code{parallel_in}?

\xpart{c} What are the register's final contents after all four shifts, and why?

\note[Hint]You need no real hardware to check Exercise~\ref{c6:ex:sipotrace} and
\ref{c6:ex:pisotrace} by hand. Work through them exactly as you would trace a truth table: one clock
edge, one row, at a time.

\exercise{\code{sipo8} and \code{piso8}}{VHDL}
\label{c6:ex:shiftregs}
Implement the chapter's two useful shift register forms in VHDL.

\xpart{a}\parttitle{\code{sipo8}.} Write an entity named \code{sipo8}, a serial-in, parallel-out
register.

\begin{booktable}{@{}p{6.6em}p{3.4em}p{11em}L@{}}
  \thead{Port} & \thead{Dir.} & \thead{Type} & \thead{Meaning} \\
  \midrule
  \code{clock} & in & \code{std_logic} & System clock. \\
  \code{reset_s2_n} & in & \code{std_logic} & Active-low, \textbf{already synchronised}.
    Asynchronous. Clears the whole register, so \code{parallel_out} reads \code{"00000000"} while it
    is held low. \\
  \code{serial_in} & in & \code{std_logic} & The incoming bit. \\
  \code{parallel_out} & out & \code{std_logic_vector(7 downto 0)} & The register's contents. \\
\end{booktable}

Follow the SIPO idiom from \secref{c6:sec:sipo}.

\modulefig[0.68]{lectures/L06/appendix/images/sipo8.png}

\xpart{b}\parttitle{\code{piso8}.} Write an entity named \code{piso8}, a parallel-in, serial-out
register.

\begin{booktable}{@{}p{6.6em}p{3.4em}p{11em}L@{}}
  \thead{Port} & \thead{Dir.} & \thead{Type} & \thead{Meaning} \\
  \midrule
  \code{clock} & in & \code{std_logic} & System clock. \\
  \code{reset_s2_n} & in & \code{std_logic} & Active-low, \textbf{already synchronised}.
    Asynchronous. Clears the whole register, so \code{serial_out} reads \code{'0'} while it is held
    low. \\
  \code{load} & in & \code{std_logic} & Capture \code{parallel_in} at the next rising edge. \\
  \code{shift} & in & \code{std_logic} & Shift one position at the next rising edge. \\
  \code{serial_in} & in & \code{std_logic} & The bit fed in at bit 0 on a shift. \\
  \code{parallel_in} & in & \code{std_logic_vector(7 downto 0)} & The value captured on a load. \\
  \code{serial_out} & out & \code{std_logic} & The register's top bit. \\
\end{booktable}

Follow the PISO idiom from \secref{c6:sec:piso}: at a load edge (\code{load = '1'}), capture
\code{parallel_in}; at a shift edge (\code{shift = '1'}), move every bit one position towards the
MSB and feed \code{serial_in} in at bit 0; drive \code{serial_out} from bit 7, so that a loaded byte
leaves most significant bit first; and let \code{load} take precedence over \code{shift} when both
are high.

\modulefig[0.68]{lectures/L06/appendix/images/piso8.png}

\xpart{c} For each of the four shift register forms (SISO, SIPO, PISO, PIPO), name a task it is the
natural choice for. Follow the summary table in \secref{c6:sec:shift}.

Derive each idiom afresh from the chapter rather than copying \code{serial_rx8}'s shift line.

\begin{selfcheck}
\textbf{Self-check.} \code{sipo8}: the ports \code{clock}, \code{reset_s2_n}, \code{serial_in} (in)
and \code{parallel_out} (\code{std_logic_vector(7 downto 0)}, out), declared in that order;
testbench in \file{exercises/sipo8/}. \code{piso8}: the ports \code{clock}, \code{reset_s2_n},
\code{load}, \code{shift}, \code{serial_in}, \code{parallel_in}
(\code{std_logic_vector(7 downto 0)}) and \code{serial_out} (out), declared in that order; testbench
in \file{exercises/piso8/}.

\textbf{Keep \code{piso8}.} \Chapref{c8:ch}'s transmitter capstone does not instantiate it, but it is
worth rereading there for the idiom's sake, since that machine shifts the byte out inline.
\end{selfcheck}

% ------------------------------------------------------------------------------------------
\exerciseset{Reading the serial receiver}

The worked example comes with a reference testbench. Run it from the example's own directory before
you start:

\begin{shell}
cd lectures/L06/serial_rx8
ghdl -a --std=93 serial_rx8.vhd serial_rx8_tb.vhd
ghdl -e --std=93 serial_rx8_tb
ghdl -r --std=93 serial_rx8_tb --assert-level=error --stop-time=10ms
\end{shell}

\exercise{The receiver's behaviour, and three real bugs}{Simulation}
\label{c6:ex:rx8}
Account for the receiver's behaviour, with reference to \secref{c6:sec:rx8}.

\xpart{a} The bit stream \code{1, 0, 1, 1, 0, 0, 1, 0} arrives, one bit per enabled clock edge. What
is on \code{data_out} during the cycle where \code{data_ready} pulses? Which arriving bit ends up in
\code{data_out(7)}, and which in \code{data_out(0)}?

\xpart{b} \code{shift_enable} goes low for three clock cycles after the fourth bit, and then high
again for the remaining four. Does the receiver still produce the same byte? Explain what the bit
counter does during the gap. What would go wrong instead if \code{shift_enable} gated only the shift
register and not the counter?

\xpart{c} Make each of the following changes in your own copy of \file{serial_rx8.vhd}, one at a
time, and \textbf{predict} what the reference testbench will report before you run it. Then run it
and account for the result:
\begin{itemize}
  \item Change the shift line to the mirrored direction,
    \code{shift_reg <= serial_in & shift_reg(7 downto 1);}.
  \item Remove the unconditional \code{data_ready <= '0';} at the top of the clocked branch.
  \item Change the bit counter's comparison from \code{7} to \code{6}.
\end{itemize}

Each of them is a real bug that someone has shipped. For each, say which of the testbench's checks
catches it, and what the failure would have looked like on a real serial link if nothing had caught
it.
